\chapter*{Abstract}

Mission planning and astronaut/operator training for planetary exploration increasingly rely on immersive VR reconstructions of candidate landing sites, but the only data available for most such sites is nadir orbital imagery (e.g.\ NASA's HiRISE, 25\,cm/px), not the ground-level, perspective views that make a VR scene feel navigable. This dissertation addresses that nadir-to-perspective domain gap in the context of ESA's ExoMars Rosalind Franklin rover mission, proposing a three-stage generative pipeline --- diffusion-based super-resolution, physics-constrained unpaired view translation, and Gaussian Splatting scene reconstruction --- to turn orbital imagery into explorable, physically plausible VR environments.

To date, the literature review and methodology have been completed, surveying the super-resolution, image-to-image translation, neural rendering, and hallucination-safety literature relevant to each stage, and formalising the proposed Neural Schr\"{o}dinger Bridge formulation, evaluation framework (FID, shadow-angle error, hallucination rate, VR frame rate), and baseline comparisons (MARSGAN, DiffusionSat, CycleGAN, UNSB-no-physics). A CycleGAN baseline for Stage 2 has been implemented, trained, and evaluated end-to-end: a real-resolution HiRISE data pipeline was built after an initial corpus was found to use browse-quality (\textasciitilde3\,m/px) imagery rather than the true 25\,cm/px RDR products, and a corrected \textasciitilde30,000-patch corpus was assembled from seven Oxia Planum observations. A 50-epoch training run on this corrected corpus was completed and evaluated (FID 269.9), and a specific, evidenced diagnosis was reached: one discriminator collapses from epoch 3 onward, most likely because the grainy, true-resolution HiRISE patches hand it an easy non-semantic shortcut over the smoother rover imagery, starving the corresponding generator of a useful gradient. In parallel, a CTX/HiRISE co-registration pipeline for the Stage 1 super-resolution corpus is under development.

This draft covers the work completed to date; addressing the CycleGAN discriminator imbalance, building the physics-constrained UNSB stage, and evaluating the full pipeline remain the subject of ongoing work.
